\section{The slopes of the period maps}\label{sec:period-slopes}

Fix the weight set $T$, its weights and the shift $u_0$ at a mesh point of
Hypothesis~\ref{hyp:one-shift-arithmetic}. Let $N_u$ be the source
of \eqref{eq:loss-source} and $C=R^\vee$ the common target. Put
\[
 d_u=\operatorname{rk}N_u,\qquad \Delta=\sum_{u\in W}d_u,\qquad
 D=\operatorname{rk}C,\qquad s=\langle\mu(N_u)-h_u\rangle .
\]
Inequality~\eqref{eq:acct-mean} asks how far $\mu(C)$ can fall below $s$.
At each slope threshold we count the target directions missed by the
period maps, then integrate the count.

\subsection{The integral maps}

Use the doubled slopes of \S\ref{sec:accounting}. For a hermitian
lattice $L$ over the coefficient field $F$, write $\lambda_i(L)$ for
its Harder--Narasimhan slopes, repeated by rank. Their mean is
$\mu(L)$. The saturated sublattice $L_{\ge a}$ has rank
$\#\{i:\lambda_i(L)\ge a\}$. Its nonzero quotients by saturated
sublattices have slope at least $a$, and every nonzero sublattice of
$L/L_{\ge a}$ has slope strictly below $a$
\citep[\S5]{gaudron-pentes}. Give sublattices and quotients the
induced and quotient norms, respectively. For a nonzero rational map
$B$ put
\[
 \operatorname{ht}(B)=[F:\mathbb Q]^{-1}
   \sum_{\sigma:F\hookrightarrow\mathbb C}\log\|B_\sigma\|_{\rm op}^{\,2}.
\]

\begin{hypothesis}[\openinput]\label{hyp:period-maps}
For every word $u$ the composite $B_u=P_uK_{m,u}\cdots K_{1,u}$ of
\S\ref{sec:theta-chain} is defined over $F$. It satisfies
$B_u(N_u)\subseteq C$ and $\operatorname{ht}(B_u)\le h_u$ when
$B_u\ne0$, with the metrics and $h_u$ of \eqref{eq:acct-height}.
\end{hypothesis}

By \eqref{eq:loss-source} the integral inclusion is equivalent to
\[
 B_u(\Lambda_{uw})\subseteq
       \pi_w^{\lfloor e_wz_{uw}\rfloor}C_w
       \qquad\text{at every finite place }w.
\]
At the inert primes this tests the final pairing of
Hypothesis~\ref{hyp:chain-endpoints} on the rescaled source.

\begin{lemma}\label{lem:period-filtration}
Under Hypothesis~\ref{hyp:period-maps},
$B_u((N_u)_{\ge a+h_u})\subseteq C_{\ge a}$ for every real $a$.
\end{lemma}

\begin{proof}
Project the indicated source to $C/C_{\ge a}$ by $B_u$. If its image
$I$ is nonzero, the quotient $L$ of the source by the kernel has
$\mu(L)\ge a+h_u$. The map $L\to I$ has operator norm at most
$\|B_{u,\sigma}\|_{\rm op}$ at each embedding. Taking determinants gives
$\mu(I)\ge\mu(L)-\operatorname{ht}(B_u)\ge a$, contradicting the
slope bound for sublattices of $C/C_{\ge a}$.
\end{proof}

\subsection{The rank estimate}

At one complex embedding $v$ of $F$, let
$A_u:\mathbb C^{d_u}\to\mathbb C^{D'}$ be matrices in the standard
hermitian norms, with $D'\ge D\ge1$. To represent the period maps they
must satisfy
\begin{equation}\label{eq:period-rows}
 \rho A_u=B_{u,v}J_u \qquad (u\in W)
\end{equation}
for one surjection $\rho:\mathbb C^{D'}\to C_v$ and isomorphisms
$J_u:\mathbb C^{d_u}\to (N_u)_v$. Here $L_v=L\otimes_{\mathcal O_F,v}
\mathbb C$. Put $\delta=1-D/D'$ and $Q=\sum_u A_uA_u^*$.
For $0\le e\le1/4$ and $\eta\ge0$ consider the bounds
\begin{equation}\label{eq:period-gram}
 \sum_u\|A_u^*A_u-I_{d_u}\|_{\rm HS}^2\le e\Delta,\qquad
 \operatorname{tr}(Q^2)\le(1+\eta)\frac{(\operatorname{tr}Q)^2}{D'} .
\end{equation}
The square of the Hilbert--Schmidt norm is $\|A\|_{\rm HS}^2
=\operatorname{tr}(A^*A)$.

\begin{lemma}\label{lem:period-rank}
Assume \eqref{eq:period-gram}. For subspaces
$V_u\subseteq\mathbb C^{d_u}$ put $p=\Delta^{-1}\sum_u\operatorname{codim}V_u$.
Then
\[
 \frac1{D'}\operatorname{codim}\sum_u A_uV_u
       \le \eta+4\sqrt e+4p .
\]
\end{lemma}

\begin{proof}
Let $P$ project orthogonally onto $(\sum_u A_uV_u)^\perp$ and put
$y=\operatorname{tr}P/D'$, $t=\operatorname{tr}Q$. By Cauchy--Schwarz,
$|t-\Delta|\le\Delta\sqrt e$, so $t\ge\Delta/2$. Since $PA_u$ kills
$V_u$, the same trace estimate on the orthogonal complements gives
\[
 \operatorname{tr}(PQ)=\sum_u\|PA_u\|_{\rm HS}^2
 \le \sum_u\|A_u|_{V_u^\perp}\|_{\rm HS}^2
 \le \Delta(p+\sqrt{pe}).
\]
On the other hand $\|Q-tI/D'\|_{\rm HS}^2\le\eta t^2/D'$ by
\eqref{eq:period-gram}. Hence $\operatorname{tr}(PQ)/t\ge
y-\sqrt{\eta y}$. Combining the bounds and using
$2\sqrt{\eta y}\le y+\eta$ and $p\le1$ gives the claim.
\end{proof}

\subsection{The mean slope}

Set $U=\max_u(\mu_{\max}(N_u)-h_u)-s\ge0$ and take
$H\ge(s+\mu_{\max}(R))_+$. By duality $\mu_{\min}(C)=-\mu_{\max}(R)$.

\begin{theorem}\label{thm:period-mean}
Under Hypothesis~\ref{hyp:period-maps} and
\eqref{eq:period-rows}--\eqref{eq:period-gram}, inequality
\eqref{eq:acct-mean} holds with
\begin{equation}\label{eq:period-error}
 e_3=\frac{(\eta+4\sqrt e)H+4\min\{U,H\}}{1-\delta}.
\end{equation}
\end{theorem}

\begin{proof}
Write $a_{ui}=\lambda_i(N_u)-h_u$ and
$p(t)=\Delta^{-1}\#\{(u,i):a_{ui}<s-t\}$. Apply
Lemma~\ref{lem:period-rank} to the inverse images under $J_u$ of
$((N_u)_{\ge s-t+h_u})_v$. By Lemma~\ref{lem:period-filtration}
and \eqref{eq:period-rows}, $\rho$ sends their combined image into
$(C_{\ge s-t})_v$. A surjection can only decrease codimension, so
\[
 \frac{\#\{i:\lambda_i(C)<s-t\}}D
     \le\frac{\eta+4\sqrt e+4p(t)}{1-\delta}.
\]
The mean of the $a_{ui}$ is $s$. Thus
\[
 \int_0^H p(t)\,dt
 \le\frac1\Delta\sum_{u,i}(s-a_{ui})_+
 =\frac1\Delta\sum_{u,i}(a_{ui}-s)_+\le U .
\]
The integral is also at most $H$. Every target slope is at least $s-H$.
Integrating the target count therefore bounds
$D^{-1}\sum_i(s-\lambda_i(C))_+$ by \eqref{eq:period-error}.
Since $s-\mu(C)$ is at most this mean, the result follows.
\end{proof}

\begin{proposition}[Word-dependent denominator costs]
\label{prop:period-denominator-average}
Assume the hypotheses of Theorem~\ref{thm:period-mean}, except for
the integral inclusions in Hypothesis~\ref{hyp:period-maps}.
Suppose positive integers $D_u$ satisfy $D_uB_u(N_u)\subseteq C$.
Keep $s,U,H$ computed with the original height allowances, and put
\[
 c_u=2\log D_u,\qquad
 \bar c=\frac{1}{\Delta}\sum_u d_uc_u,\qquad
 c_{\min}=\min_u c_u.
\]
Then
\begin{equation}\label{eq:period-denominator-average}
 s\le\mu(C)+\bar c+
 \frac{(\eta+4\sqrt e)H+
       4\min\{U+\bar c-c_{\min},H\}}{1-\delta}.
\end{equation}
In particular, if $\delta\le1/2$ and
$(\eta+4\sqrt e)H\le6b$, then
\[
 s\le\mu(C)+8U+12b+9\bar c-8c_{\min}
       \le\mu(C)+8U+12b+9\bar c.
\]
For a common denominator $D_u=D_x$, the first bound recovers the
original mean-slope error plus $2\log D_x$.
\end{proposition}

\begin{proof}
Use auxiliary maps $\widetilde B_u=D_uB_u$ and height allowances
$\widetilde h_u=h_u+c_u$. At every embedding multiplication by
$D_u$ increases the squared logarithmic operator norm by $c_u$,
so these are valid height bounds; zero maps remain zero. The new
mean is $\widetilde s=s-\bar c$. Its spread satisfies
\[
 \widetilde U
 =\max_u\bigl(\mu_{\max}(N_u)-h_u-c_u\bigr)-s+\bar c
 \le U+\bar c-c_{\min}.
\]
Since $\bar c\ge0$, the old target cutoff $H$ still works.
Set $\widetilde J_u=D_u^{-1}J_u$. These remain isomorphisms and
$\rho A_u=\widetilde B_{u,v}\widetilde J_u$, so the matrices and
their Gram bounds are unchanged. The proof of
Theorem~\ref{thm:period-mean} uses only the integral inclusions,
the height bounds, the row identity and the Gram estimates, so it
applies to these auxiliary data. Translate back by adding $\bar c$.
This proves \eqref{eq:period-denominator-average}.
For the displayed consequence use
$(1-\delta)^{-1}\le2$ and $\bar c-c_{\min}\ge0$.
If the $c_u$ are all equal, $\widetilde U=U$ and the only added
term is that common value. No invariance of the deficit proportions
is asserted when the denominators vary with $u$.
\end{proof}

The same argument applies to the scalar-row formulation of
Corollary~\ref{cor:bessel-mean}. Keep $I_u$ and replace
$\kappa_u$ by $\kappa_u/D_u$ in
$\rho A_u=\kappa_u B_{u,v}I_u$. The filtration inclusion and
the unchanged Gram bounds give the same estimate, also when
$\kappa_u=0$.

The rescaling in \S\ref{sec:inert-gain} also bounds $U$. At $w\mid p$
we have $\Lambda_{uw}\subseteq p^{\tau_p}M_{u,w}$, with
$\tau_p=0$ off $\mathcal P$. Hence $N_u\subseteq\mathfrak b_u M_u$ for
the fractional ideal $\mathfrak b_u$ with orders
$e_w\tau_p-\lfloor e_wz_{uw}\rfloor$. Twisting each sublattice by
$\mathfrak b_u$ adds $-2\sum_w\ell_w\operatorname{ord}_w(\mathfrak b_u)$
to its slope. With the rounding term $R_u$ of Theorem~\ref{thm:loss-source},
\[
 \mu_{\max}(N_u)\le\mu_{\max}(M_u)+\Gamma+\mathcal L-E_u-R_u .
\]
Subtract the average of \eqref{eq:loss-second-error} after deducting
$h_u$ from each source slope. The common $\Gamma+\mathcal L$ cancels,
giving, under the hypotheses of Theorem~\ref{thm:loss-source},
\begin{equation}\label{eq:period-spread}
 U\le \max_u\bigl(\mu_{\max}(M_u)-E_u-h_u-R_u\bigr)
        -\langle\mu^M-E-h\rangle+\langle e_{2,u}\rangle .
\end{equation}

\begin{hypothesis}[\openinput]\label{hyp:period-frame}
At each required mesh point the normalized Bessel period columns of the
same words, restricted to the prescribed type at $T_u$, realize
\eqref{eq:period-rows} at one embedding. Their Gram matrices satisfy
\eqref{eq:period-gram}. The bounds for $e,\eta,\delta,H$ and the source
spread in \eqref{eq:period-spread} put \eqref{eq:period-error} within
the $e_3$ allowance of \S\ref{sec:arithmetic-bound}, uniformly on the
chosen windows.
\end{hypothesis}

For the Bessel periods of \citep[\S13]{adriazola-grh}, the second
bound in \eqref{eq:period-gram} uses all pairs, since
\[
 \operatorname{tr}(Q^2)=\sum_{u,u'\in W}\|A_u^*A_{u'}\|_{\rm HS}^2,
 \qquad \operatorname{tr}Q=\sum_{u\in W}\|A_u\|_{\rm HS}^2 .
\]
Under Hypotheses~\ref{hyp:period-maps} and~\ref{hyp:period-frame},
Theorem~\ref{thm:period-mean} gives \eqref{eq:acct-mean} at the chosen
shift. In particular, $D\ge D'/2$ and $(\eta+4\sqrt e)H\le6b$ give
$e_3\le8U+12b$.
